% Manual de uso de la plantilla UTI para trabajos de titulación
% Compilar desde la raíz del proyecto con:
%   latexmk -pdf -outdir=manual manual/Manual_de_Uso.tex
\documentclass[12pt,a4paper,oneside]{report}

\usepackage[T1]{fontenc}
\usepackage[utf8]{inputenc}
\usepackage[spanish,es-tabla]{babel}
\usepackage{lmodern}
\usepackage{geometry}
\geometry{top=3cm,bottom=3cm,left=4cm,right=3cm}
\usepackage{setspace}
\onehalfspacing
\usepackage{array}
\usepackage{longtable}
\usepackage{tabularx}
\usepackage{xcolor}
\usepackage{graphicx}
\usepackage{listings}
\usepackage{enumitem}
\usepackage{amssymb}
\usepackage{fancyhdr}
\usepackage[hidelinks]{hyperref}

\definecolor{UTIMorado}{HTML}{422568}
\definecolor{UTINaranja}{HTML}{F56600}
\definecolor{CodigoFondo}{HTML}{F4F2F7}

\hypersetup{
  pdftitle={Manual de Uso de la Plantilla para Trabajos de Integración Curricular},
  pdfauthor={Carrera de Ingeniería en Tecnologías de la Información},
  colorlinks=true,
  linkcolor=UTIMorado,
  urlcolor=UTIMorado,
  citecolor=UTIMorado
}

\pagestyle{fancy}
\fancyhf{}
\fancyhead[L]{\small Manual de Uso}
\fancyhead[R]{\small Ingeniería en Tecnologías de la Información}
\fancyfoot[C]{\thepage}
\renewcommand{\headrulewidth}{0.4pt}
\renewcommand{\footrulewidth}{0pt}
\setlength{\headheight}{14pt}

\fancypagestyle{plain}{%
  \fancyhf{}%
  \fancyfoot[C]{\thepage}%
  \renewcommand{\headrulewidth}{0pt}%
}

\setlength{\parindent}{1.25cm}
\setlength{\parskip}{0pt}
\setlength{\emergencystretch}{2em}
\setlist{itemsep=0.25em,topsep=0.4em}

\lstset{
  basicstyle=\ttfamily\small,
  backgroundcolor=\color{CodigoFondo},
  frame=single,
  rulecolor=\color{UTIMorado},
  breaklines=true,
  columns=fullflexible,
  showstringspaces=false,
  keepspaces=true,
  upquote=true,
  literate=
    {á}{{\'a}}1 {é}{{\'e}}1 {í}{{\'i}}1 {ó}{{\'o}}1 {ú}{{\'u}}1
    {Á}{{\'A}}1 {É}{{\'E}}1 {Í}{{\'I}}1 {Ó}{{\'O}}1 {Ú}{{\'U}}1
    {ñ}{{\~n}}1 {Ñ}{{\~N}}1 {ü}{{\"u}}1 {¿}{{?`}}1
}

\newcommand{\advertencia}[1]{%
  \begin{center}
    \fcolorbox{UTINaranja}{CodigoFondo}{%
      \begin{minipage}{0.88\textwidth}
        \textbf{Importante:} #1
      \end{minipage}%
    }
  \end{center}
}

\begin{document}
\hypersetup{pageanchor=false}

\begin{titlepage}
  \thispagestyle{empty}
  \begin{center}
    \vspace*{1.5cm}
    {\color{UTIMorado}\rule{\textwidth}{2pt}}\\[1.2cm]
    {\Large\bfseries UNIVERSIDAD INDOAMÉRICA\par}
    \vspace{0.5cm}
    {\large\bfseries FACULTAD DE INGENIERÍAS\par}
    \vspace{0.3cm}
    {\large\bfseries CARRERA DE INGENIERÍA EN TECNOLOGÍAS DE LA INFORMACIÓN\par}
    \vfill
    {\Huge\bfseries\color{UTIMorado} MANUAL DE USO\par}
    \vspace{0.5cm}
    {\LARGE Plantilla para Trabajos de Integración Curricular\par}
    \vspace{1cm}
    {\large Guía práctica para estudiantes y docentes tutores\par}
    \vfill
    \begin{tabular}{rl}
      \textbf{Código documental:} & UTI--TIC--MU--01 \\
      \textbf{Versión:} & 1.0 \\
      \textbf{Fecha:} & 20 de julio de 2026 \\
    \end{tabular}
    \vfill
    {\color{UTINaranja}\rule{\textwidth}{2pt}}
  \end{center}
\end{titlepage}

\hypersetup{pageanchor=true}
\pagenumbering{roman}

\chapter*{Cómo utilizar este manual}
\addcontentsline{toc}{chapter}{Cómo utilizar este manual}

Este documento guía la preparación completa del Trabajo de Integración Curricular usando la plantilla LaTeX institucional. Se recomienda seguir los capítulos en orden durante la configuración inicial y consultar las secciones técnicas cuando se incorporen tablas, figuras, ecuaciones, referencias o anexos.

\advertencia{No modifique \texttt{uti\_tfm.sty}. Ese archivo contiene las reglas institucionales controladas por el Comité Curricular. Los cambios ordinarios se realizan en \texttt{main.tex}, \texttt{Preliminares/}, \texttt{Capítulos/}, \texttt{Anexos/} y \texttt{referencias.bib}.}

\tableofcontents
\clearpage
\pagenumbering{arabic}

\chapter{Preparación del entorno}

\section{Requisitos}

Para trabajar con la plantilla se necesita:

\begin{itemize}
  \item una distribución LaTeX actualizada, como TeX Live o MiKTeX;
  \item un editor compatible, por ejemplo TeXstudio, TeXmaker, Overleaf o Visual Studio Code con LaTeX Workshop;
  \item BibTeX para procesar las referencias IEEE;
  \item un visor PDF;
  \item la fuente Century Gothic si se utilizará XeLaTeX o LuaLaTeX.
\end{itemize}

Si Century Gothic no está disponible, puede compilarse con pdfLaTeX. La plantilla utilizará una fuente alternativa de aspecto similar.

\section{Copia de trabajo}

Antes de editar:

\begin{enumerate}
  \item duplique la carpeta completa del proyecto;
  \item asigne a la copia un nombre identificable;
  \item mantenga la estructura de carpetas;
  \item no cambie nombres que contengan tildes sin actualizar también los comandos \texttt{\textbackslash input};
  \item utilice control de versiones o copias de respaldo periódicas.
\end{enumerate}

\section{Archivos que debe editar}

\begin{longtable}{|p{4cm}|p{8.6cm}|}
  \hline
  \textbf{Archivo o carpeta} & \textbf{Uso} \\
  \hline
  \texttt{main.tex} & Datos generales, autoría, tutor, título inglés, línea de investigación y orden del documento. \\
  \hline
  \texttt{Preliminares/} & Textos legales, dedicatoria, agradecimiento, resumen y abstract. \\
  \hline
  \texttt{Capítulos/} & Contenido académico y técnico principal. \\
  \hline
  \texttt{Anexos/} & Evidencias y material complementario. \\
  \hline
  \texttt{referencias.bib} & Datos bibliográficos de todas las fuentes citadas. \\
  \hline
  \texttt{Capítulos/fig/} & Diagramas, gráficos y figuras. \\
  \hline
  \texttt{Capítulos/img/} & Fotografías y capturas de pantalla. \\
  \hline
\end{longtable}

\chapter{Configuración de datos generales}

\section{Título inglés}

En \texttt{main.tex}, reemplace el texto de ejemplo:

\begin{lstlisting}
\UTIDatosTituloIngles{ENGLISH PROJECT TITLE}
\end{lstlisting}

Este dato se utiliza únicamente en el abstract. Debe ser una traducción académica fiel del título español.

\section{Selección de línea y sublínea}

El comando recibe dos índices:

\begin{lstlisting}
\UTISeleccionInvestigacion{linea}{sublinea}
\end{lstlisting}

Por ejemplo, para seleccionar ``Gestión de la información'' y ``Análisis y procesamiento de datos'' se usa:

\begin{lstlisting}
\UTISeleccionInvestigacion{3}{2}
\end{lstlisting}

Mientras el tutor no haya validado la selección, mantenga:

\begin{lstlisting}
\UTISeleccionInvestigacion{0}{0}
\end{lstlisting}

\subsection{Catálogo de selección}

\begin{longtable}{|p{2cm}|p{10.6cm}|}
  \hline
  \textbf{Índices} & \textbf{Línea y sublínea} \\
  \hline
  \multicolumn{2}{|p{12.6cm}|}{\textbf{Línea 1: Análisis, diseño e implementación de infraestructura tecnológica}} \\
  \hline
  \texttt{\{1\}\{1\}} & Diseño y optimización de redes de comunicación. \\
  \hline
  \texttt{\{1\}\{2\}} & Ciberseguridad y protección de infraestructuras tecnológicas. \\
  \hline
  \texttt{\{1\}\{3\}} & Sistemas de almacenamiento y computación en la nube. \\
  \hline
  \texttt{\{1\}\{4\}} & Automatización y gestión de infraestructuras inteligentes. \\
  \hline
  \multicolumn{2}{|p{12.6cm}|}{\textbf{Línea 2: Redes informáticas de comunicación}} \\
  \hline
  \texttt{\{2\}\{1\}} & Arquitecturas de redes avanzadas. \\
  \hline
  \texttt{\{2\}\{2\}} & Protocolos y estándares de comunicación. \\
  \hline
  \texttt{\{2\}\{3\}} & Ciberseguridad en redes. \\
  \hline
  \texttt{\{2\}\{4\}} & Redes inalámbricas e Internet de las Cosas (IoT). \\
  \hline
  \multicolumn{2}{|p{12.6cm}|}{\textbf{Línea 3: Gestión de la información}} \\
  \hline
  \texttt{\{3\}\{1\}} & Almacenamiento y recuperación de información. \\
  \hline
  \texttt{\{3\}\{2\}} & Análisis y procesamiento de datos. \\
  \hline
  \texttt{\{3\}\{3\}} & Gobernanza y seguridad de la información. \\
  \hline
  \texttt{\{3\}\{4\}} & Ciclo de vida de la información. \\
  \hline
  \texttt{\{3\}\{5\}} & Sistemas de gestión del conocimiento. \\
  \hline
  \texttt{\{3\}\{6\}} & Visualización de información. \\
  \hline
  \texttt{\{3\}\{7\}} & Tecnologías emergentes en la gestión de información. \\
  \hline
  \texttt{\{3\}\{8\}} & Gestión de datos no estructurados. \\
  \hline
  \multicolumn{2}{|p{12.6cm}|}{\textbf{Línea 4: Diseño, desarrollo e integración de sistemas basados en arquitecturas web, móviles y sistemas embebidos}} \\
  \hline
  \texttt{\{4\}\{1\}} & Diseño y desarrollo de aplicaciones web. \\
  \hline
  \texttt{\{4\}\{2\}} & Desarrollo de aplicaciones móviles. \\
  \hline
  \texttt{\{4\}\{3\}} & Integración de sistemas embebidos. \\
  \hline
  \texttt{\{4\}\{4\}} & Interoperabilidad entre plataformas. \\
  \hline
  \texttt{\{4\}\{5\}} & Seguridad en sistemas web, móviles y embebidos. \\
  \hline
  \texttt{\{4\}\{6\}} & Optimización del rendimiento. \\
  \hline
  \texttt{\{4\}\{7\}} & Aplicación de tecnologías emergentes. \\
  \hline
  \texttt{\{4\}\{8\}} & Sistemas basados en la nube y edge computing. \\
  \hline
\end{longtable}

La línea y sublínea seleccionadas aparecen automáticamente en la tabla de Área de Estudio del capítulo II.

\section{Portada}

El bloque principal tiene la siguiente estructura:

\begin{lstlisting}
\UTIPortada
{LOGO-UTI-1.png}
{FACULTAD DE INGENIERÍAS}
{INGENIERÍA EN TECNOLOGÍAS DE LA INFORMACIÓN}
{TÍTULO DEL PROYECTO}
{Trabajo de Titulación modalidad Trabajo de Integración Curricular}
{\UTIDatosAutores{Nombres Apellidos}{femenino}}
{\UTIDatosTutor{Nombres Apellidos}{masculino}}
{Quito, Ecuador}
{2026}
\end{lstlisting}

Cambie únicamente el contenido dentro de cada par de llaves. No elimine argumentos ni modifique su orden.

\section{Uno o dos autores}

Para un autor:

\begin{lstlisting}
{\UTIDatosAutores{Ana Perez}{femenino}}
\end{lstlisting}

Para dos autores:

\begin{lstlisting}
{\UTIDatosAutores{Ana Perez}{femenino}
  [Carlos Lopez][masculino]}
\end{lstlisting}

Valores permitidos para el género: \texttt{masculino} y \texttt{femenino}. Estos valores controlan pronombres, etiquetas, conjugaciones y títulos profesionales en las páginas preliminares.

\section{Tutor o tutora}

\begin{lstlisting}
{\UTIDatosTutor{Mgtr. Nombre Apellido}{masculino}}
% o
{\UTIDatosTutor{Mgtr. Nombre Apellido}{femenino}}
\end{lstlisting}

Incluya únicamente el grado académico más alto que corresponda mostrar.

\chapter{Edición de páginas preliminares}

\section{Documentos legales}

Revise los archivos \texttt{01} a \texttt{04} de \texttt{Preliminares/}. Los nombres, el título, la ciudad, el año y el género se toman automáticamente de \texttt{main.tex}. Debe completar manualmente:

\begin{itemize}
  \item fecha exacta;
  \item cédula, dirección, correo y teléfono;
  \item nombres de lectores o miembros del tribunal;
  \item firmas cuando corresponda en la versión final.
\end{itemize}

No cambie el contenido legal sin autorización institucional.

\section{Dedicatoria y agradecimiento}

Edite \texttt{05\_Dedicatoria.tex} y \texttt{06\_Agradecimiento.tex}. Si existen dos autores, la plantilla genera una página para cada persona. Reemplace las instrucciones de ejemplo por el texto definitivo.

\section{Resumen ejecutivo}

Edite el tercer argumento de \texttt{\textbackslash UTIResumenPagina} en \texttt{08\_Resumen.tex}. Requisitos:

\begin{itemize}
  \item un solo párrafo;
  \item máximo 300 palabras;
  \item problema, objetivo, metodología, resultados y conclusiones;
  \item redacción impersonal, clara y precisa;
  \item de tres a cuatro descriptores, ordenados alfabéticamente.
\end{itemize}

Elimine \texttt{\textbackslash textit} y los corchetes de instrucción al introducir el texto final.

\section{Abstract}

Edite \texttt{09\_Abstract.tex}. Debe ser equivalente al resumen, no una traducción automática sin revisión. La cabecera se presenta en inglés y el título se obtiene de \texttt{\textbackslash UTIDatosTituloIngles}.

\chapter{Redacción del cuerpo}

\section{Regla general}

Todo texto en cursiva y entre corchetes es una instrucción. Debe eliminarse y sustituirse por contenido académico definitivo.

Para añadir apartados utilice:

\begin{lstlisting}
\section{Nombre de la sección}
Texto del apartado.

\subsection{Nombre de la subsección}
Texto de la subsección.
\end{lstlisting}

No escriba manualmente números delante de los títulos. La plantilla controla su presentación.

Las secciones aparecerán automáticamente como \textbf{1.1}, \textbf{1.2}, \textbf{2.1}, etc.; las subsecciones como \textbf{1.1.1}. En el Índice de Contenidos, los capítulos se muestran explícitamente como \textbf{CAPÍTULO I}, \textbf{CAPÍTULO II}, etc.

\section{Capítulo I: Introducción}

Desarrolle contextualización, estado del arte, problema, justificación, objetivos, alcance y fundamentación teórica. El estado del arte debe comparar fuentes pertinentes y relacionarlas con el proyecto, no limitarse a una lista de resúmenes.

\section{Capítulo II: Metodología}

Complete el diagnóstico, el Área de Estudio, la metodología, las técnicas de recolección, los requerimientos, el diseño y la justificación técnica.

En la tabla de Área de Estudio, la línea, sublínea y campo son automáticos. Edite directamente dominio, área, aspectos, objeto y periodo de estudio.

\section{Capítulo III: Implementación, Pruebas y Resultados}

Documente la solución implementada, las actividades ejecutadas, las pruebas y los resultados medibles. Compare la situación inicial con la posterior y conserve evidencias verificables.

\section{Capítulo IV: Conclusiones y Recomendaciones}

Cada conclusión debe relacionarse con los objetivos y resultados. Las recomendaciones deben derivarse de hallazgos reales. Los trabajos futuros deben plantear extensiones factibles.

\chapter{Inserción de elementos}

\section{Figuras}

Utilice figuras para diagramas, arquitecturas, gráficos o esquemas:

\begin{lstlisting}
Como se observa en la Figura~\ref{fig:arquitectura}, ...

\begin{figure}[htbp]
  \centering
  \includegraphics[width=0.80\textwidth]
    {Capítulos/fig/arquitectura.png}
  \caption{Arquitectura de la solución.}
  \label{fig:arquitectura}
\end{figure}
\elaboradopor{Estudiante}
\fuente{Elaboración propia}
\end{lstlisting}

La etiqueta debe ser única. No escriba ``Figura 1'' manualmente; use \texttt{\textbackslash ref}.

\section{Imágenes y capturas}

Para fotografías o capturas utilice el entorno \texttt{imagen}:

\begin{lstlisting}
\begin{imagen}[htbp]
  \centering
  \includegraphics[width=0.80\textwidth]
    {Capítulos/img/panel.png}
  \caption{Panel de administración implementado.}
  \label{img:panel}
\end{imagen}
\elaboradopor{Estudiante}
\fuente{Captura del entorno de pruebas}
\end{lstlisting}

\section{Tablas}

\begin{lstlisting}
\begin{table}[htbp]
  \centering
  \caption{Resultados de las pruebas.}
  \label{tab:resultados}
  \begin{tabular}{|l|c|c|}
    \hline
    Metrica & Antes & Despues \\
    \hline
    Tiempo de respuesta & 850 ms & 320 ms \\
    \hline
  \end{tabular}
\end{table}
\elaboradopor{Estudiante}
\fuente{Resultados experimentales}
\end{lstlisting}

Si la tabla excede el ancho disponible, reduzca el contenido, use columnas de párrafo o divídala. No reduzca el tamaño hasta volverla ilegible.

\section{Ecuaciones}

\begin{lstlisting}
\begin{equation}
\label{eq:disponibilidad}
Disponibilidad\,(\%) =
\frac{T_{operativo}}{T_{total}} \times 100
\end{equation}
\addeqtolist{Cálculo de disponibilidad del servicio}
\end{lstlisting}

Mencione la ecuación mediante \texttt{Ecuación\textasciitilde\textbackslash ref\{eq:disponibilidad\}}.

\section{Referencias cruzadas}

Ejemplos:

\begin{lstlisting}
Capítulo~\ref{chap:metodologia}
Tabla~\ref{tab:resultados}
Figura~\ref{fig:arquitectura}
Imagen~\ref{img:panel}
Ecuación~\ref{eq:disponibilidad}
\end{lstlisting}

El símbolo \texttt{\textasciitilde} evita que el nombre y el número se separen al final de una línea.

\chapter{Citas y referencias IEEE}

\section{Agregar una fuente}

Incluya cada fuente en \texttt{referencias.bib}. Ejemplo:

\begin{lstlisting}
@article{apellido2026titulo,
  author  = {Nombre Apellido and Nombre Apellido},
  title   = {Título del artículo},
  journal = {Nombre de la revista},
  volume  = {10},
  number  = {2},
  pages   = {10--20},
  year    = {2026},
  doi     = {10.0000/ejemplo}
}
\end{lstlisting}

La clave \texttt{apellido2026titulo} debe ser única.

\section{Citar dentro del texto}

\begin{lstlisting}
La arquitectura propuesta reduce la latencia del servicio \cite{apellido2026titulo}.

Distintos estudios sustentan esta decisión \cite{clave1,clave2,clave3}.
\end{lstlisting}

No escriba manualmente números como \texttt{[1]}. BibTeX asigna y actualiza la numeración.

\section{Calidad de fuentes}

Priorice artículos científicos, libros académicos, normas técnicas, documentación oficial y actas de congresos. Verifique autor, título, año, editorial o revista, páginas, DOI y URL cuando corresponda.

\chapter{Glosario y anexos}

\section{Glosario}

Edite \texttt{Capítulos/06\_Glosario.tex}:

\begin{lstlisting}
\begin{description}
  \item[API] Interfaz de programación de aplicaciones.
  \item[IoT] Internet de las Cosas.
\end{description}
\end{lstlisting}

Incluya únicamente términos relevantes que requieran aclaración.

\section{Anexos}

El primer anexo se crea así:

\begin{lstlisting}
\Anexo{Manual técnico de la solución}

Contenido del anexo.
\end{lstlisting}

Para añadir otro anexo:

\begin{lstlisting}
\Anexo{Instrumento de entrevista}

Contenido del segundo anexo.
\end{lstlisting}

La plantilla asigna automáticamente A, B, C y actualiza el índice de anexos.

\chapter{Compilación del documento}

\section{Método recomendado}

Desde la carpeta principal del proyecto:

\begin{lstlisting}
latexmk -xelatex main.tex
\end{lstlisting}

Este método requiere Century Gothic. Si la fuente no está instalada:

\begin{lstlisting}
latexmk -pdf main.tex
\end{lstlisting}

\section{Compilación manual}

Con XeLaTeX:

\begin{lstlisting}
xelatex main.tex
bibtex main
xelatex main.tex
xelatex main.tex
\end{lstlisting}

Con pdfLaTeX:

\begin{lstlisting}
pdflatex main.tex
bibtex main
pdflatex main.tex
pdflatex main.tex
\end{lstlisting}

La bibliografía actual utiliza BibTeX. No ejecute biber salvo que el sistema bibliográfico sea cambiado oficialmente.

\section{Cuándo recompilar}

Compile nuevamente cuando cambie citas, bibliografía, títulos, etiquetas, índices, anexos, tablas, figuras o imágenes. Dos pasadas finales eliminan la mayoría de referencias pendientes.

\chapter{Solución de problemas}

\begingroup
\small
\singlespacing
\begin{longtable}{|p{4.2cm}|p{4.2cm}|p{4.2cm}|}
  \hline
  \textbf{Problema} & \textbf{Causa probable} & \textbf{Solución} \\
  \hline
  Century Gothic no encontrada & La fuente no está instalada. & Instálela conforme a la licencia o compile con pdfLaTeX. \\
  \hline
  Aparecen signos \texttt{??} & Faltan pasadas de compilación o la etiqueta no existe. & Compile dos veces y revise \texttt{\textbackslash label}. \\
  \hline
  No aparece una referencia & BibTeX no se ejecutó o la clave es incorrecta. & Ejecute la secuencia completa y compare la clave con \texttt{referencias.bib}. \\
  \hline
  Imagen no encontrada & Ruta, nombre, extensión o mayúsculas incorrectas. & Verifique el archivo y copie la ruta exacta. \\
  \hline
  Código de línea inválido & El par seleccionado no existe. & Consulte el catálogo y corrija ambos índices. \\
  \hline
  Tabla fuera del margen & Columnas demasiado anchas. & Use columnas \texttt{p\{...\}}, reduzca texto o divida la tabla. \\
  \hline
  Índice desactualizado & Archivos auxiliares pendientes. & Compile hasta que no se soliciten nuevas ejecuciones. \\
  \hline
  Texto legal en género incorrecto & Género mal configurado. & Use exactamente \texttt{masculino} o \texttt{femenino}. \\
  \hline
  El abstract muestra el título español & No se configuró el título inglés. & Complete el título inglés en \texttt{main.tex}. \\
  \hline
\end{longtable}
\endgroup

\chapter{Lista de verificación final}

\section{Contenido}

\begin{itemize}
  \item[$\square$] Se eliminaron todos los textos de instrucción entre corchetes.
  \item[$\square$] Los objetivos coinciden con las conclusiones y resultados.
  \item[$\square$] Todas las figuras, tablas, imágenes y ecuaciones se mencionan en el texto.
  \item[$\square$] Cada cita tiene una referencia y cada referencia fue citada.
  \item[$\square$] El resumen y el abstract son equivalentes y no superan 300 palabras.
  \item[$\square$] La línea y sublínea fueron validadas por el tutor.
  \item[$\square$] El periodo y objeto de estudio están claramente definidos.
\end{itemize}

\section{Formato}

\begin{itemize}
  \item[$\square$] La portada contiene datos definitivos.
  \item[$\square$] Los márgenes y la numeración de páginas son correctos.
  \item[$\square$] Los índices están actualizados.
  \item[$\square$] Las leyendas y fuentes son legibles.
  \item[$\square$] No existen referencias con \texttt{??}.
  \item[$\square$] El PDF abre correctamente y todas las páginas son visibles.
  \item[$\square$] El nombre del archivo final cumple la convención indicada por la carrera.
\end{itemize}

\section{Entrega}

Antes de entregar, conserve:

\begin{itemize}
  \item PDF final;
  \item carpeta completa del proyecto LaTeX;
  \item archivos de imágenes y diagramas originales;
  \item base \texttt{referencias.bib};
  \item respaldos o repositorio de control de versiones;
  \item anexos y evidencias solicitadas por tutor o tribunal.
\end{itemize}

\advertencia{La aceptación del formato no reemplaza la revisión académica. El estudiante debe atender las observaciones del tutor, lectores, tribunal y Comité Curricular.}

\end{document}